\chapter{Internal Complex Structure and Projective Geometry of \(2\)-Level States}
\label{chap:geometria-proyectiva-bloch}
\index{internal root of minus one}
\index{Paley!conference matrix}
\index{Witt design}
\index{Bloch sphere}
\index{Möbius!action}
\index{binary icosahedral group}

The phase geometry employed by the quantum realization has an algebraic genealogy predating complex notation. Its point of departure is the finite Paley--Witt incidence, from which an endomorphism whose square is minus the identity is obtained. Realification preserves that operator as the phase generator, and projectivization of a two-level space leads to the Bloch sphere. On that sphere, the spinorial action of \(\mathrm{SU}(2)\), its expression through Möbius transformations, and the binary icosahedral group form successive realizations of the same structure.

The causal chain developed in this chapter is
\begin{equation}
 \boxed{
 \begin{gathered}
 C_P\longrightarrow A_W^2=-I_6,\\
 A_W^2=-I_6\longrightarrow
 \mathbb F_3^6\simeq\mathbb F_9^3
 \longrightarrow (E,J_E),\\
 (E,J_E)\longrightarrow
 \mathbb P(\mathcal H_J)\simeq\mathbb{CP}^{1}
 \longrightarrow S^2,\\
 S^2\longrightarrow
 \bigl\{\mathrm{SU}(2),\ \mathop{\text{M\"ob}}(S^2),\ 2A_5\bigr\}.
 \end{gathered}}
 \label{eq:cadena-paley-bloch-anexo}
\end{equation}
Each arrow declares its domain. The finite quadratic structure, the
orthogonal realification, and the projective geometry remain related, with
their scalar fields explicitly separated.

\section{The incidence of Paley as an algebraic origin of the phase}

Let \(C_P\in M_6(\mathbb Z)\) be the symmetric conference matrix
\begin{equation}
 C_P=
 \begin{pmatrix}
 0& 1& 1& 1& 1& 1\\
 1& 0& 1&-1&-1& 1\\
 1& 1& 0& 1&-1&-1\\
 1&-1& 1& 0& 1&-1\\
 1&-1&-1& 1& 0& 1\\
 1& 1&-1&-1& 1& 0
 \end{pmatrix}.
 \label{eq:paley-conferencia-anexo}
\end{equation}
Their rows satisfy
\begin{equation}
 C_P C_P^{\mathsf T}=5I_6.
 \label{eq:paley-ortogonalidad-anexo}
\end{equation}
By reducing modulo three we obtain
\begin{equation}
 A_W:=C_P\bmod 3\in M_6(\mathbb F_3).
 \label{eq:paley-reduccion-anexo}
\end{equation}

\begin{theorem}[Internal root of minus the identity]
\label{thm:raiz-interna-menos-uno-anexo}
The endomorphism \(J_W:v\mapsto vA_W\) of
\(V_W=\mathbb F_3^6\) satisfies
\begin{equation}
 J_W^2=-I_{V_W}.
 \label{eq:Jw-cuadrado-anexo}
\end{equation}
Consequently, \(V_W\) admits a canonical module structure over
\(\mathbb F_3[X]/(X^2+1)\) once the Paley matrix has been fixed.
\end{theorem}

\begin{proof}
The matrix \(C_P\) is symmetric. Reducing
\eqref{eq:paley-ortogonalidad-anexo} modulo three yields
\[
 A_W^2=5I_6\bmod3=2I_6=-I_6.
\]
Therefore, the evaluation \(X\mapsto J_W\) annihilates the ideal
\((X^2+1)\) and defines the action of the indicated quotient.
\end{proof}

The polynomial \(X^2+1\) is irreducible over \(\mathbb F_3\), since its
values in \(0,1,2\) are \(1,2,2\). Therefore,
\begin{equation}
 \mathbb F_9:=\mathbb F_3[\iota]/(\iota^2+1)
 \label{eq:F9-desde-paley-anexo}
\end{equation}
is a field, and \(J_W\) realizes multiplication by
\(\iota\).

\begin{corollary}[Descent from six ternary coordinates to three quadratic coordinates]
\label{cor:F3seis-F9tres-anexo}
The rule
\begin{equation}
 (a+b\iota)v:=av+bJ_Wv,
 \qquad a,b\in\mathbb F_3,
 \label{eq:accion-F9-anexo}
\end{equation}
converts \(V_W\) into a vector space of dimension three over
\(\mathbb F_9\). In the coordinates fixed by \(A_W\),
\begin{equation}
 \mathbb F_3^6\simeq\mathbb F_9^3.
 \label{eq:F3seis-igual-F9tres-anexo}
\end{equation}
\end{corollary}

\begin{proof}
The identity \(J_W^2=-I\) verifies the multiplication rules of
\(\mathbb F_9\). Since
\([\mathbb F_9:\mathbb F_3]=2\) and
\(\dim_{\mathbb F_3}V_W=6\), the dimension over \(\mathbb F_9\) is three.
\end{proof}

The same operator organizes the ternary code via the graph.
\begin{equation}
 C_W=\{(q,qA_W):q\in\mathbb F_3^6\}\subset\mathbb F_3^{12}.
 \label{eq:codigo-grafo-J-anexo}
\end{equation}
If \(G=(I_6\mid A_W)\), then
\[
 GG^{\mathsf T}=I_6+A_W A_W^{\mathsf T}=0.
\]
Thus, the internal square root of minus the identity participates simultaneously in the
quadratic structure and the isotropy of the code. This dual function explains
why phase and exceptional incidence share an antecedent while
preserving distinct readers.

\section{Orthogonal realification of the phase}

The analytic phase is constructed on a real space before a complex
chart is chosen. Let \((E,g)\) be a real Euclidean space of even dimension, and let
\(J_E\in\operatorname{End}(E)\) be an operator such that
\begin{equation}
 J_E^2=-I,
 \qquad
 g(J_Eu,J_Ev)=g(u,v).
 \label{eq:estructura-compleja-ortogonal-anexo}
\end{equation}
The second identity implies \(J_E^{\mathsf T}=-J_E\).

\begin{proposition}[Phase as an orthogonal group of one parameter]
\label{prop:fase-ortogonal-J-anexo}
For every \(\theta\in\mathbb R\),
\begin{equation}
 R_J(\theta):=e^{\theta J_E}
 =\cos\theta\,I+\sin\theta\,J_E
 \label{eq:rotacion-J-anexo}
\end{equation}
is orthogonal and satisfies
\begin{equation}
 R_J(\theta)R_J(\phi)=R_J(\theta+\phi),
 \qquad R_J(2\pi)=I.
 \label{eq:grupo-fase-J-anexo}
\end{equation}
\end{proposition}

\begin{proof}
The exponential series separates into even and odd powers and, using
\(J_E^2=-I\), produces the trigonometric expression. The antisymmetry of
\(J_E\) gives
\(R_J(\theta)^{\mathsf T}=R_J(-\theta)\); therefore,
\(R_J(\theta)^{\mathsf T}R_J(\theta)=I\). The composition law follows from the fact
that all factors are functions of the same operator.
\end{proof}

The choice of the chart \(J_E\leftrightarrow i\) turns \(E\) into a
complex vector space: for \(a,b\in\mathbb R\),
\begin{equation}
 (a+ib)u:=au+bJ_Eu.
 \label{eq:carta-compleja-J-anexo}
\end{equation}
In that chart, \(R_J(\theta)\) is written as multiplication by
\(e^{i\theta}\). The complex notation publishes the phase already transported by the real
operator.

\begin{remark}[Two scales, the same quadratic relation]
The operator \(J_W\) acts on \(\mathbb F_3^6\); the operator
\(J_E\) acts on a real inner-product space. The prolongation
\(J_W\rightsquigarrow J_E\) is a realization that preserves the relation
\(J^2=-I\) and additionally specifies a topology, norm, and completion.
This order avoids using the continuous analytic structure as a premise of
the finite incidence.
\end{remark}

\section{The complex projective line and the Bloch sphere}

Let \(\mathcal H_J\) be a Hilbert space of complex dimension two obtained
through the chart \eqref{eq:carta-compleja-J-anexo}. A normalized pure state
is a vector
\begin{equation}
 \psi=\begin{pmatrix}z_0\\ z_1\end{pmatrix},
 \qquad |z_0|^2+|z_1|^2=1,
 \label{eq:espinor-normalizado-anexo}
\end{equation}
and its projective physical state is the ray
\([\psi]=\{e^{i\vartheta}\psi:\vartheta\in\mathbb R\}\).
The ray space is \(\mathbb P(\mathcal H_J)\simeq\mathbb{CP}^{1}\).

We define the coordinates
\begin{equation}
 \begin{aligned}
 r_1&=2\operatorname{Re}(\overline z_0z_1),\\
 r_2&=2\operatorname{Im}(\overline z_0z_1),\\
 r_3&=|z_0|^2-|z_1|^2.
 \end{aligned}
 \label{eq:coordenadas-bloch-anexo}
\end{equation}

The spherical representation of two-level states is formulated in the
coordinates introduced by Bloch \cite{Bloch1946}.

\begin{theorem}[Projective Bloch Identification]
\label{thm:CP1-esfera-bloch-anexo}
The map
\begin{equation}
 \mathcal B:\mathbb{CP}^{1}\longrightarrow S^2,
 \qquad [\psi]\longmapsto(r_1,r_2,r_3),
 \label{eq:mapa-bloch-anexo}
\end{equation}
is a smooth bijection. The rank-one projector associated with \([\psi]\) is
\begin{equation}
 \rho_\psi=|\psi\rangle\langle\psi|
 =\frac12\bigl(I_2+r_1\sigma_1+r_2\sigma_2+r_3\sigma_3\bigr),
 \label{eq:proyector-bloch-anexo}
\end{equation}
where \(\sigma_1,\sigma_2,\sigma_3\) are the Pauli matrices.
\end{theorem}

\begin{proof}
Multiplication of \(\psi\) by a common phase leaves the three
coordinates invariant. Moreover,
\[
 r_1^2+r_2^2+r_3^2
 =4|z_0|^2|z_1|^2+(|z_0|^2-|z_1|^2)^2=1.
\]
Every \(\boldsymbol r\in S^2\) determines the positive projector \(\rho=(I+\boldsymbol r\cdot\boldsymbol\sigma)/2\), which
satisfies \(\rho^2=\rho\) and \(\operatorname{tr}\rho=1\); its image is a single ray. The
direct matrix formula for \(|\psi\rangle\langle\psi|\) yields \eqref{eq:proyector-bloch-anexo}. The usual affine
charts of \(\mathbb{CP}^{1}\) make the map and its inverse smooth.
\end{proof}

The Bloch sphere records the projective class and preserves, through
\(\rho_\psi\), all expectation values of two-level Hermitian observables.
The global phase belongs to the fibre \(S^1\) of the Hopf bundle
\cite{Hopf1931}
\begin{equation}
 S^1\longrightarrow S^3\longrightarrow S^2.
 \label{eq:hopf-bloch-anexo}
\end{equation}
The HMT distinction between enriched state and published coordinate finds
a precise realization here: the normalized vector belongs to \(S^3\), the
projective reader publishes a point of \(S^2\), and the fiber preserves the phase that
the projection identifies.

\section{Spinorial Action, Rotations, and Möbius Transformations}

Every element of \(\mathrm{SU}(2)\) can be written as
\begin{equation}
 U=
 \begin{pmatrix}
 \alpha&\beta\\
 -\overline\beta&\overline\alpha
 \end{pmatrix},
 \qquad |\alpha|^2+|\beta|^2=1.
 \label{eq:SU2-parametrizacion-anexo}
\end{equation}
Acts linearly on spinors and, by conjugation, on projectors.

\begin{theorem}[Double covering and action on the sphere]
\label{thm:SU2-SO3-bloch-anexo}
for each \(U\in\mathrm{SU}(2)\) there exists a unique
\(R_U\in\mathrm{SO}(3)\) such that
\begin{equation}
 U(\boldsymbol r\cdot\boldsymbol\sigma)U^\dagger
 =(R_U\boldsymbol r)\cdot\boldsymbol\sigma.
 \label{eq:accion-SU2-bloch-anexo}
\end{equation}
The map \(U\mapsto R_U\) is a surjective homomorphism with kernel
\(\{\pm I_2\}\).
\end{theorem}

\begin{proof}
Conjugation preserves the three-dimensional real space of trace-free Hermitian
matrices and its inner product
\(\langle A,B\rangle=\tfrac12\operatorname{tr}(AB)\). It therefore induces an
orthogonal transformation of \(\mathbb R^3\). Continuity and the value at the
identity fix a positive determinant. If conjugation is trivial on the
three Pauli matrices, \(U\) commutes with every complex matrix \(2\times2\),
and by Schur's lemma it is scalar; the determinant-one condition leaves
\(U=\pm I_2\). The axis--angle representation shows that every rotation has
two opposite preimages.
\end{proof}

In the affine chart \(z=z_0/z_1\) of \(\mathbb{CP}^{1}\), the same action takes
the form
\begin{equation}
 z\longmapsto
 \frac{\alpha z+\beta}
      {-\overline\beta z+\overline\alpha}.
 \label{eq:mobius-SU2-anexo}
\end{equation}
The Bloch rotation and the fractional transformation are therefore two
coordinates of the same projective action. The former employs the Hermitian
projector; the latter, the Riemann sphere.

\section{The icosahedral specialization and the binary group \texorpdfstring{\(2A_5\)}{2A5}}

The grupo rotacional of the icosaedro is a subgrupo of
\(\mathrm{SO}(3)\) isomorfo to \(A_5\). Its preimagen by the recubrimiento of the
\cref{thm:SU2-SO3-bloch-anexo} defines the grupo icosaedrico binario
\begin{equation}
 2A_5:=\pi^{-1}(A_5)\subset\mathrm{SU}(2).
 \label{eq:def-2A5-anexo}
\end{equation}

\begin{proposition}[Icosahedral central extension]
\label{prop:extension-central-2A5-anexo}
There exists an exact sequence
\begin{equation}
 1\longrightarrow\{\pm I_2\}
 \longrightarrow2A_5
 \longrightarrow A_5
 \longrightarrow1,
 \label{eq:sucesion-2A5-anexo}
\end{equation}
and \(|2A_5|=120\). An icosahedral rotation of orders three or five lifts to
elements of orders six or ten, respectively.
\end{proposition}

\begin{proof}
The restriction of \(\pi:\mathrm{SU}(2)\to\mathrm{SO}(3)\) to the preimage of \(A_5\) is surjective and preserves the kernel \(\{\pm I_2\}\). Exactness and the order theorem yield \(|2A_5|=2|A_5|=120\). A rotation through angle \(2\pi/n\) is represented in \(\mathrm{SU}(2)\) with half-angle \(\pi/n\); after \(n\) iterations it reaches \(-I_2\) and after \(2n\), the identity.
\end{proof}

The icosahedral specialization establishes a rigorous hinge between the geometry of two-level states, Möbius transformations, and the finite-symmetry branch. The development of McKay theory and of the exceptional structures subsequently uses the representation of \(2A_5\) and its modules; its additional content is developed in the specific theorems of that branch.

\section{Functorial composition \texorpdfstring{\(\mathbb F_3^6\to\mathbb F_9^3\to\mathbb{CP}^{1}\simeq S^2\to\mathrm{SU}(2)\to\mathrm{SO}(3)\to 2A_5\)}{F3^6 -> F9^3 -> CP1 -> S2 -> SU(2) -> SO(3) -> 2A5}}

The relation \(J^2=-I\) traverses four categories while preserving its type:
\begin{enumerate}
 \item in \(\mathbb F_3^6\), defines the quadratic extension
       \(\mathbb F_9\) and organizes the Witt incidence;
 \item in a real orthogonal space, generates the phase flow
       \(e^{\theta J_E}\);
 \item in complex dimension two, projectivization produces
       \(\mathbb{CP}^{1}\simeq S^2\) and the Bloch projectors;
 \item the action of \(\mathrm{SU}(2)\) descends to rotations and
       Möbius transformations, and its icosahedral preimage produces
       \(2A_5\).
\end{enumerate}
The internal complex structure thus has a finite provenance and an explicit
analytic realization. The Bloch sphere publishes the geometry of rays; HMT
genealogy additionally preserves the route, orientation, carry, and memory
that precede that publication.
